% GENERATED FILE. Edit canonical lessons, then run npm run generate:books.
% canonical-source-hash: fnv1a64:db8d648b2d366185
% canonical-lessons: LA-C92-mus, LA-C92-gallina, LA-C92-asinus, LA-C92-elephantus, LA-C92-cervus

\chapter{Mouse, Hen, Donkey, Elephant, Deer}
\label{ch:la-mouse-hen-donkey-elephant-deer}
\hlchaptermodality{92}

\begin{chapteropening}
\textbf{By the end of this chapter:} \emph{I can name a mouse, a hen, a donkey, an elephant and a deer.}

Complete the last lesson of chapter 92: I can name a mouse, a hen, a donkey, an elephant and a deer.
\end{chapteropening}

\section[mūs, mūris]{mūs, mūris — a mouse}

\label{lesson:LA-C92-mus}



\begin{quote}
Before the new one: say the Latin for a bear, then the Latin for a snake.
\end{quote}

\subsection*{You'll want to know: mūs, mūris}
\textbf{mūs, mūris} — \textquotedblleft{}a mouse\textquotedblright{}.

\begin{cousinweb}
\textbf{mūs} is inherited, usually traced to Proto-Indo-European \emph{*muh2s}, \textquotedblleft{}mouse\textquotedblright{} (EtymDB 2.1). Greek \emph{mys}, \textquotedblleft{}mouse\textquotedblright{} (Lewis \& Short), and English \emph{mouse} are cousins.
\end{cousinweb}

\begin{grammarlens}[title={What you've built}]
One more word for this chapter.
\end{grammarlens}

\subsection*{Your turn}
\begin{itemize}
\raggedright
  \item \emph{Say it:} \emph{mūs}
  \item \emph{Say it:} \emph{mūs}, once more
  \item \emph{Say it:} say \emph{mūs}
  \item \emph{Recall:} say the Latin for a bear, then the Latin for a snake, then say \emph{mūs} again
\end{itemize}

\begin{tcolorbox}[breakable,colback=teal!4,colframe=teal!35!black,title={Before you move on}]
What does mūs, mūris mean? (\textquotedblleft{}A mouse\textquotedblright{}.) Say it once more.
\end{tcolorbox}
\section[gallīna, gallīnae]{gallīna, gallīnae — a hen}

\label{lesson:LA-C92-gallina}



\begin{quote}
Before the new one: say the Latin for a snake, then the Latin for a mouse.
\end{quote}

\subsection*{You'll want to know: gallīna, gallīnae}
\textbf{gallīna, gallīnae} — \textquotedblleft{}a hen\textquotedblright{}. It lays the \textbf{ōvum}.

\begin{cousinweb}
\textbf{gallīna} is built on \emph{gallus}, \textquotedblleft{}a rooster\textquotedblright{} (Lewis \& Short): the hen is the rooster's female. Where \emph{gallus} comes from is uncertain.
\end{cousinweb}

\begin{grammarlens}[title={What you've built}]
One more word for this chapter.
\end{grammarlens}

\subsection*{Your turn}
\begin{itemize}
\raggedright
  \item \emph{Say it:} \emph{gallīna}
  \item \emph{Say it:} \emph{gallīna}, once more
  \item \emph{Say it:} say \emph{gallīna}
  \item \emph{Recall:} say the Latin for a snake, then the Latin for a mouse, then say \emph{gallīna} again
\end{itemize}

\begin{tcolorbox}[breakable,colback=teal!4,colframe=teal!35!black,title={Before you move on}]
What does gallīna, gallīnae mean? (\textquotedblleft{}A hen\textquotedblright{}.) Say it once more.
\end{tcolorbox}
\section[asinus, asinī]{asinus, asinī — a donkey}

\label{lesson:LA-C92-asinus}



\begin{quote}
Before the new one: say the Latin for a mouse, then the Latin for a hen.
\end{quote}

\subsection*{You'll want to know: asinus, asinī}
\textbf{asinus, asinī} — \textquotedblleft{}a donkey\textquotedblright{}.

\begin{cousinweb}
\textbf{asinus} is probably a loan from an eastern language, as Lewis \& Short already suspected; Greek \emph{onos}, \textquotedblleft{}donkey\textquotedblright{}, may share the source, which is unknown. English \emph{ass} comes from \emph{asinus}.
\end{cousinweb}

\begin{grammarlens}[title={What you've built}]
One more word for this chapter.
\end{grammarlens}

\subsection*{Your turn}
\begin{itemize}
\raggedright
  \item \emph{Say it:} \emph{asinus}
  \item \emph{Say it:} \emph{asinus}, once more
  \item \emph{Say it:} say \emph{asinus}
  \item \emph{Recall:} say the Latin for a mouse, then the Latin for a hen, then say \emph{asinus} again
\end{itemize}

\begin{tcolorbox}[breakable,colback=teal!4,colframe=teal!35!black,title={Before you move on}]
What does asinus, asinī mean? (\textquotedblleft{}A donkey\textquotedblright{}.) Say it once more.
\end{tcolorbox}
\section[elephantus, elephantī]{elephantus, elephantī — an elephant}

\label{lesson:LA-C92-elephantus}



\begin{quote}
Before the new one: say the Latin for a hen, then the Latin for a donkey.
\end{quote}

\subsection*{You'll want to know: elephantus, elephantī}
\textbf{elephantus, elephantī} — \textquotedblleft{}an elephant\textquotedblright{}.

\begin{cousinweb}
\textbf{elephantus} was borrowed from Greek \emph{elephas}, \emph{elephantos}, \textquotedblleft{}elephant, ivory\textquotedblright{} (Lewis \& Short). Where Greek got the word is uncertain.
\end{cousinweb}

\begin{grammarlens}[title={What you've built}]
One more word for this chapter.
\end{grammarlens}

\subsection*{Your turn}
\begin{itemize}
\raggedright
  \item \emph{Say it:} \emph{elephantus}
  \item \emph{Say it:} \emph{elephantus}, once more
  \item \emph{Say it:} say \emph{elephantus}
  \item \emph{Recall:} say the Latin for a hen, then the Latin for a donkey, then say \emph{elephantus} again
\end{itemize}

\begin{tcolorbox}[breakable,colback=teal!4,colframe=teal!35!black,title={Before you move on}]
What does elephantus, elephantī mean? (\textquotedblleft{}An elephant\textquotedblright{}.) Say it once more.
\end{tcolorbox}
\section[cervus, cervī]{cervus, cervī — a deer, a stag}

\label{lesson:LA-C92-cervus}



\begin{quote}
Before the new one: say the Latin for a donkey, then the Latin for an elephant.
\end{quote}

\subsection*{You'll want to know: cervus, cervī}
\textbf{cervus, cervī} — \textquotedblleft{}a deer, a stag\textquotedblright{}.

\begin{cousinweb}
\textbf{cervus} is inherited, usually traced to Proto-Indo-European \emph{*krh2wos}, \textquotedblleft{}deer, stag\textquotedblright{}, from \emph{*kerh2-}, \textquotedblleft{}horn\textquotedblright{} (EtymDB 2.1): the horned one. \emph{cornū}, \textquotedblleft{}horn\textquotedblright{}, comes from the same root (Lewis \& Short), and English \emph{hart} is a cousin.
\end{cousinweb}

\begin{grammarlens}[title={What you've built}]
One more word for this chapter.
\end{grammarlens}

\subsection*{Your turn}
\begin{itemize}
\raggedright
  \item \emph{Say it:} \emph{cervus}
  \item \emph{Say it:} \emph{cervus}, once more
  \item \emph{Say it:} say \emph{cervus}
  \item \emph{Recall:} say the Latin for a donkey, then the Latin for an elephant, then say \emph{cervus} again
\end{itemize}

\begin{tcolorbox}[breakable,colback=teal!4,colframe=teal!35!black,title={Before you move on}]
What does cervus, cervī mean? (\textquotedblleft{}A deer, a stag\textquotedblright{}.) Say it once more.
\end{tcolorbox}
